% Source-derived chapter 03: Deformation-theoretic preliminaries
% Original source: canonical-representations-surface-groups
\chapter{Deformation-theoretic preliminaries}
\label{canonical-representations-surface-groups-chapter-03}
\source{canonical-representations-surface-groups}

\begin{remark}[Source section]
\label{canonical-representations-surface-groups-chapter-03-source}
\source{canonical-representations-surface-groups}
\source{canonical-representations-surface-groups}
This chapter reproduces the corresponding section of the retrieved
LaTeX source, including its definitions, statements, examples, and
proofs. It has not yet been translated into Lean.
\notready
\end{remark}

% The source section title was: Deformation-theoretic preliminaries
\label{section:deformation-theory}
\source{canonical-representations-surface-groups}
In this section we introduce some deformation theory of group representations. In \autoref{section:main-proof}, we will use the theory developed here in conjunction with the vanishing results proven in \autoref{section:vanishing} to show that MCG-finite representations have certain rigidity properties.
\subsection{Deformations of representations}\label{subsection:deformation-basics}
\source{canonical-representations-surface-groups}
We begin by recalling a standard description of deformations of representations in terms of cohomological data.

Recall 
\autoref{notation:adjoint}:
for $R$ a ring and a representation $\rho: G\to \on{GL}_r(R)$, we use $\on{ad}(\rho)$
to denote the adjoint $G$-representation obtained by composing $\rho$ with
$\on{GL}_r(R)\to \on{GL}(\mathfrak{pgl}_r(R)).$ 

Fix a group $G$ and a representation $\rho_0: G\to GL_r(\mathbb{C})$. Let $\on{Art}$ be the category of local Artin $\mathbb{C}$-algebras with residue field $\mathbb{C}$. If $A\in\on{Art}$, with
maximal ideal $\mathfrak{m}_A$, we say that a representation $\rho_A: G\to
\on{GL}_r(A)$ is a deformation of $\rho_0$ if $\rho_A=\rho_0 \bmod \mathfrak{m}_A$. The \emph{constant} deformation is the deformation obtained via the composition $$G\overset{\rho_0}{\to}GL_r(\mathbb{C})\to GL_r(A).$$ We say a representation $\rho: G\to GL_r(A)$ is \emph{conjugate to a constant representation} if there exists a matrix $M\in GL_r(A)$ such that $M\rho M^{-1}$ factors through $GL_r(\mathbb{C})$.
Say a deformation $\rho_A$ of $\rho_0$ has \emph{constant determinant} if $\det\rho_A=\det\rho_0$, regarded as a map $G\to A^\times$ via the natural inclusion $\mathbb{C}^\times\hookrightarrow A^\times$.

We now define deformations of representations over square-zero extensions in
$\on{Art}$.
  Given a square-zero extension $$0\to I\to B\to A\to 0$$ with $A, B\in\on{Art}$, and a representation $\rho_A: G\to GL_r(A)$, a deformation of $\rho_A$ is a representation $\rho_B: G\to \on{GL}_r(B)$ with $\rho_B=\rho_A\bmod I$. Two such deformations $\rho_B, \rho_B'$ are \emph{equivalent} if there exists a matrix $M\in \on{GL}_r(B)$ with $\rho_B=M\rho_B'M^{-1}$, such that $M=\on{id}\bmod I$. 
  \begin{lemma}
\source{canonical-representations-surface-groups}
\notready
	\label{lemma:tangent-space}
\source{canonical-representations-surface-groups}
	Fix a group $G$ and a representation $\rho_0 : G \to \on{GL}_r(\mathbb C)$. Consider $$0\to \epsilon A\to A[\epsilon]/\epsilon^2\to A\to 0$$ with $A\in\on{Art}$, and let $\rho_A: G\to \on{GL}_r(A)$ be a deformation of $\rho_0$ over $A$ with constant determinant.
	The set of equivalence classes of deformations of $\rho_A$ to an $A[\epsilon]/\epsilon^2$-representation with constant determinant is naturally in bijection with $H^1(G, \on{ad}(\rho_A))$.
%	is isomorphic to the $\mathbb C$-vector space of
%	conjugacy classes of deformations 
%	$\rho: G \to \on{GL}_r(\mathbb C[\varepsilon]/(\varepsilon^2))$
%	of $\rho_0$ whose associated determinant map $\det \rho :
%	G \to \mathbb G_m(\mathbb C[\varepsilon]/(\varepsilon^2))$ is
%	constant, i.e., $\det \rho = \det \rho_0 \otimes_{\mathbb C} \mathbb C[\varepsilon]/(\varepsilon^2)$. 
%Here ``conjugacy classes of deformations," means two deformations are the same if they are conjugate 
%by a matrix of the form $I_r+\epsilon A, A\in \mathfrak{gl}_r(\mathbb{C})$.
	Moreover, this description is functorial in $G$.
\end{lemma}
\begin{proof}
See \cite[Chapter V, \S21, Proposition 1]{mazur1997introduction}.
%See also \cite[Propositions 4.4 and
%4.6]{klevdalP:g-rigid-local-systems-are-integral} in the irreducible case,
%taking $N = 0$ in 
%\cite[Proposition 4.4]{klevdalP:g-rigid-local-systems-are-integral}.
\end{proof}

\begin{lemma}
\source{canonical-representations-surface-groups}
\notready\label{lemma:split-constancy}
\source{canonical-representations-surface-groups}
	Let $$1\to N\to G\to Q\to 1$$ be a short exact sequence of groups (i.e.~$N$ is normal in $G$ and $Q=G/N$). Let $A\in\on{Art}$ and $$\rho: G\to GL_r(A)$$ be a representation, and suppose that $H^0(Q, H^1(N, \on{ad}(\rho)|_N))=0.$ If $\rho_\epsilon$ is any deformation of $\rho$ to $A[\epsilon]/\epsilon^2$ with constant determinant, then $\rho_\epsilon|_N$ is equivalent to $\rho\otimes_A A[\epsilon]/\epsilon^2|_N$.
\end{lemma}
\begin{proof}
	By \autoref{lemma:tangent-space}, the set of equivalence classes of deformations of $\rho$ to $A[\epsilon]/\epsilon^2$ is naturally in bijection with $H^1(G, \on{ad}(\rho))$, and the set of equivalence classes of deformations of $\rho|_N$ is naturally in bijection with $H^1(N, \on{ad}(\rho)|_N)$. By functoriality, the restriction map is given by the natural map $$H^1(G, \on{ad}(\rho))\to H^1(N, \on{ad}(\rho)|_N),$$ which, by the $5$-term exact sequence from the Hochschild-Serre spectral sequence, factors through $H^0(Q, H^1(N, \on{ad}(\rho)|_N))=0.$ Hence $\rho_\epsilon|_N$ is equivalent to $\rho\otimes_A A[\epsilon]/\epsilon^2|_N$ as desired.
\end{proof}

\subsection{A criterion for constancy}
The next proposition gives a criterion for a deformation of a group representation to induce the trivial deformation on a normal subgroup.
This will be used in the proof of the semisimple version of our main theorem,
\autoref{theorem:finite-image-semisimple}.

We define  $$B_n=\mathbb{C}[t]/t^{n+1},$$ $$R_n=B_{n-1}[\epsilon]/\epsilon^2=\mathbb{C}[t, \epsilon]/(t^n,\epsilon^2).$$ Let $d_n: B_n\to R_{n}$ be the map given by $$d_n: f(t)\mapsto f(t)+\epsilon f'(t),$$ where $f'(t)$ is the derivative of $f(t)$. Note that $d_n$ is injective.

Below we say that an element of $GL_r(B_n)$ is constant if it lies in $GL_r(\mathbb{C})\subset GL_r(B_n)$. We say that an element of $GL_r(\mathbb{C}[[t]])$ is constant mod $t^{n+1}$ if its image in $GL_r(B_n)$ is constant. Likewise, a representation $\rho$ into $GL_r(B_n)$ is constant if it factors through $GL_r(\mathbb{C})$; it is conjugate to a constant representation if there exists $M\in GL_r(B_n)$ such that $M\rho M^{-1}$ is constant. Similarly, a representation into $GL_r(\mathbb{C}[[t]])$ is constant mod $t^{n+1}$ if its image in $GL_r(B_n)$ is constant.
\begin{proposition}
\source{canonical-representations-surface-groups}
\notready\label{prop:stacky-split-constancy}
\source{canonical-representations-surface-groups}
	Let $$1\to N\to G\to Q\to 1$$ be a short exact sequence of groups. Let $$\rho_0: G\to GL_r(\mathbb{C})$$ be a representation, and suppose that for all $n$ and all deformations $\gamma: G\to GL_r(B_n)$ of $\rho_0$ with $\gamma|_N$ constant, we have $H^0(Q, H^1(N, \on{ad}(\gamma)|_N))=0.$ If $$\rho_\infty: G\to GL_r(\mathbb{C}[[t]])$$ is any representation with $\rho_\infty=\rho_0\bmod t$ and constant determinant, then $\rho_\infty|_N$ is conjugate to a constant representation. That is, there exists $M_\infty\in GL_r(\mathbb{C}[[t]])$ such that $M_\infty\rho_\infty|_NM_\infty^{-1}$ factors through $GL_r(\mathbb{C})$.
\end{proposition}
\begin{proof}
We wish to construct $M_\infty\in GL_r(\mathbb{C}[[t]])$ such that $M_\infty\rho_\infty|_NM_\infty^{-1}$ is constant, i.e.~factors through $GL_r(\mathbb{C})$. We do so by successive approximation. For all $m\geq 0$, set $\rho_m: G\to GL_r(B_m)$ to be $\rho_\infty \bmod t^{m+1}$.

Set $S_1=\on{id}$. Suppose we have found $S_n\in GL_r(\mathbb{C}[[t]])$ such
that $S_n\rho_\infty|_N S_n^{-1}$ is constant modulo $t^n$.  We claim it
suffices to construct an element $M_n\in GL_r(B_n), M_n=\on{id}\bmod t^n$, such
that $M_nS_n\rho_n|_N S_n^{-1}M_n^{-1}$ is constant. Indeed, let
$M_1 = \id$ and for $n \geq 1$, let
$\widetilde{M_n}$ be an arbitrary lift of $M_n$ to $GL_r(\mathbb{C}[[t]])$. 
Then, the representation $\widetilde{M_n}S_n\rho_\infty|_NS_n^{-1}\widetilde{M_n}^{-1}$ is
constant mod $t^{n+1}$. Set $S_{n+1}=\widetilde{M_n}S_n$, so that $S_{n+1} =
\widetilde{M_n} \cdot \widetilde{M_{n-1}} \cdots \widetilde{M_1}$ by induction.
Then, setting $$M_\infty:=\lim_{n\to
\infty}\widetilde{M_n}\cdot\widetilde{M_{n-1}}\cdots
\widetilde{M_1}=\lim_{n\to\infty} S_n,$$ we claim that the limit converges as
$\widetilde{M_n}\to \on{id}$, and $M_\infty\rho_\infty|_N M_{\infty}^{-1}$ is
constant. This claim holds because $M_\infty=S_n\bmod t^n$ for all $n$, so we
have $$M_\infty \rho_\infty|_N M_\infty^{-1}\bmod t^n=S_n\rho_\infty|_N
S_n^{-1}\bmod t^n,$$ which is constant for all $n$ by construction. 

We now construct the desired matrix $M_n\in GL_r(B_n)$. Replacing $\rho_\infty$ by $S_n\rho_\infty S_n^{-1}$, we may assume $\rho_\infty|_N$ is constant modulo $t^n$, i.e.~$\rho_{n-1}|_N$ is constant. We wish to find $M_n\in GL_r(B_n)$ such that
\begin{enumerate}
\item $M_n=\on{id} \bmod t^n$, and
\item $M_n\rho_n|_NM_n^{-1}$ is constant.
\end{enumerate}

Given $g\in N$, we let $\rho_n(g)'\in \on{Mat}_{r\times r}(B_{n-1})$ be the matrix obtained by differentiating the entries of $\rho_n(g)$. The representation $d_n\circ \rho_n|_N$ is given by $$g\mapsto \rho_0(g)+\epsilon\rho_n(g)',$$ as $\rho_n|_N$ is constant mod $t^{n}$.

Note that $d_n\circ \rho_n|_N: N\to GL_r(R_n)$ is constant mod $\epsilon$ (as it is equal to $\rho_{n-1}|_N$ mod $\epsilon$). By hypothesis, $$H^0(Q, H^1(N, \on{ad}(\rho_{n-1})|_N))=0,$$ and so by \autoref{lemma:split-constancy}, $d_n\circ\rho_n|_N$ is conjugate to the constant representation $$\rho_{n-1}|_N\otimes_{B_{n-1}} R_n=\rho_0|_N\otimes_{\mathbb{C}} R_n$$ 
by some matrix $$\on{id}+\epsilon\sum_{i=0}^{n-1}C_it^i\in GL_r(R_n),$$ where
the $C_i\in \on{Mat}_{r\times r}(\mathbb{C})$. The idea of the remainder of the proof is to view the matrix above as a vector field which we may flow along to make $\rho_n|_N$ constant; we find the desired conjugating matrix $M_n$ by ``integrating" this vector field.

%Note that $\rho_{n-1}|_N\otimes_{B_{n-1}} R_n$ is constant as $\rho_{n-1}|_N$ is constant, i.e. $$\rho_{n-1}|_N\otimes_{B_{n-1}} R_n=\rho_0|_N\otimes_{\mathbb{C}} R_n.$$
We compute that for $g\in N$,
\begin{align*}
	(\rho_0\otimes_{\mathbb{C}} R_n)(g) &= (\rho_{n-1}\otimes_{B_{n-1}}R_n)(g)\\
	&=(\on{id}+\epsilon\sum_{i=0}^{n-1}C_it^i)(d_n\circ \rho_n(g))(\on{id}-\epsilon\sum_{i=0}^{n-1}C_it^i)\\
	&=(\on{id}+\epsilon\sum_{i=0}^{n-1}C_it^i)(\rho_0(g)+\epsilon\rho_n(g)')(\on{id}-\epsilon\sum_{i=0}^{n-1}C_it^i)\\
&=\rho_0(g)+\epsilon\rho_n(g)'+\epsilon\sum_{i=0}^{n-1}[C_i, \rho_{0}(g)]t^i.
\end{align*}
As $\rho(g)$ is constant mod $t^n$, $\rho(g)'$ is $0$ mod $t^{n-1}$. Hence equating coefficients, we find $$[C_i, \rho_0(g)]=0 \text{ for $i<n-1$}$$ and  $$\rho_n(g)'=[\rho_0(g), C_{n-1}]t^{n-1}.$$
Now set $$M_n=\on{id}+\frac{C_{n-1}}{n}t^n\in GL_r(B_n).$$ We claim $M_n\rho_n|_N M_n^{-1}$is constant. Indeed,
\begin{align*}
M_n\rho_n(g) M_n^{-1} &= (\on{id}+\frac{C_{n-1}}{n}t^n)\rho_n(g)(\on{id}-\frac{C_{n-1}}{n}t^n)\\
&=\rho_n(g)+\frac{t}{n}[t^{n-1}C_{n-1}, \rho_n(g)]\\
&=\rho_n(g)+\frac{t}{n}[t^{n-1}C_{n-1}, \rho_0(g)]\\
&=\rho_n(g)-\frac{t}{n}\rho_n(g)'.
\end{align*}
But the only non-vanishing coefficient of $\rho(g)'$ is the coefficient of $t^{n-1}$, so the above expression is constant as desired.
\end{proof}

%As in \autoref{notation:versal-family}, let $\pi^\circ:\mathscr{C}^\circ\to
%\mathscr{M}$ be a punctured versal family.

%\begin{proposition}
%	\label{proposition:vanishing-deformations}
%	Let $\mathbb W$ be a local system on $\mathscr C^\circ$.
%	Fix a point $m \in \mathscr M$,
%and let $C^\circ := (\pi^\circ)^{-1}(m)$.
%	Suppose
%	$H^0(\mathscr M, R^1 \pi^\circ_* \mathbb W) = 0$.
%	Then the natural restriction map
%	$r: H^1(\mathscr C^\circ, \mathbb W) \to H^1(C^\circ,
%	\mathbb W|_{C^\circ})$ vanishes.
	
%	In particular, suppose $\mathbb W = \on{ad}\mathbb U$, for $\mathbb U$ a
%	local system on $\mathscr C^\circ$, and $\mathbb U_t$ is a deformation of $\mathbb U$ over a local Artin $\mathbb{C}$-algebra $A$ with residue field $\mathbb{C}$, with $\det\mathbb U_t|_{C^\circ}$
%	constant. Then $\mathbb{U}_t$
%	restricts to the trivial deformation of $\mathbb U|_{C^\circ}$ on $C^\circ$.
%\end{proposition}

%It follows from the five-term exact sequence associated to the Leray spectral
%sequence for the map $\mathscr C^\circ \to \mathscr M$ that the natural restriction map
%	$H^1(\mathscr C^\circ, \mathbb W) \to H^1(C^\circ,
%	\mathbb W|_{C^\circ})$ factors through 
%	$H^0(\mathscr M, R^1 \pi^\circ_* \mathbb W)$.
%	\begin{align*}
%		H^0(\mathscr M, R^1 \pi^\circ_* \mathbb W)\subset 
%	H^0(m, R^1 (\pi^\circ|_m)_* \mathbb W|_{C^\circ}) \simeq H^1(C^\circ_m,
%	\mathbb W|_{C^\circ_m})	
%	\end{align*}
%	Hence, the map $r$ vanishes because
%	$H^0(\mathscr M, R^1 \pi^\circ_* \mathbb W) = 0$.
%	The final statement follows from \autoref{lemma:tangent-space} by induction on the length of $A$, interpreting the map $$H^1(\mathscr{C}^\circ, \on{ad}\mathbb{U})\to H^1(C^\circ, \on{ad}\mathbb{U}|_{C^\circ})$$ as the map sending a first-order deformation of $\mathbb{U}$ to a deformation of its restriction to $C^\circ$.
%\end{proof}
